\documentclass[12pt]{article}

\usepackage[a4paper, total={6in, 8in}]{geometry}
\usepackage{textcomp}

\begin{document}
\setlength{\parindent}{0pt}

\def \t {\textrightarrow}
\def \v {\vspace{1em}}
\def \bi {\begin{itemize}}
\def \ei {\end{itemize}}
\def \s[#1] {\section*{#1}}
\def \ss[#1] {\subsection*{#1}}
\def \sss[#1] {\subsubsection*{#1}}

\s[Spiritualismo]
Movimento di pensiero che si colloca alla reazione al positivismo che si trova in europa \\
È il movimento che + si oppone, anche prima all'esistenzialismo \t pos. ha ridotto l uomo a un ente di natura \\
Positivismo ha dato una lettura incompleta, parziale \t uomo è ente di natura, ma viene contestata questa riduzione esclusiva \\
È stato letto tutta la dimensione umana come oggetto di natura \t un oggetto che poteva essere studiato con le scienze naturali \t positivismo viene accusato di essere stato assoluto nella definizione di uomo e di scienza \\
Nuell uomo c'è dell altro \t troviamo dimensione di interiorita, liberta, coscienza, che va indagato \t con dei metodi diversi pero rispetto alla scienza di natura \\
Bergson è l'esponente \t e parte dallo studio dei positivisti \t lui studia i positivisti, e indaga il positivismo \t e afferma che il positivismo non afferma la sua promessa \\
Positivismo è fedele ai fatti \t ma in realta è fedele a dei fatti che possono essere indagati con certi metodi \t li latri li ignora \\
Spiritualismo non è in scontro \t riconosce valore del pos. e sicuramente metodo scientifico serve ed è utile \t bisogna pero recuperare una dimensione che il positivismo aveva disconosciuto, perche aveva indagato con i metodi della scienza \\
Bisogna collegare determinismo e cio che appartiene all uomo \t no volonta di negazione, ma di recupero

\v

Lo spiritualismo è una reazione al positivismo, ma è anche contro all'idealismo \t l'idealismo ha identita tra finito e infinito, mentre spir. separa \t e anzi enfatizza la trascendenza dell'infinito \\
C'è salto qualitativo tra dimensione dell assoluto e il reale \t c'è l'assoluto infinito, e poi esiste l'uomo in quanto spirito finito \\
Con spirito non si intende una dimensione religiosa, ma una dimensione di liberta, interiorita, riflessione, responsabilita \t tutto cio che appartiene allo spirito, ovvero interiorita dell uomo, che era stata disconosciuta da pos.

\v

Spiritualismo si diffonde in tutta europa \t ma la culla sara la Francia \t qua c'è anche Le Chie, Tru, e poi Bergson che è l'esponente

\s[Bergson]
Lui sviluppa una filosofia che diventa il punto di riferimento del pensiero francese

\ss[Vita]
Nasce a parigi nel 59 \t studia meccanica e matematica, poi si dedica a filosofia \\
Si laurea in filosofia, insegna in diversi licei e poi fa tesi di dottorato nel 89? "Saggio sui dati immediati della coscienza" \\
La tesi viene pubblicata e ottiene enorme successo, e anche il suo lavoro del 96 ovvero "Materia e memoria" = diventa famoso, e insegna nei licei contemporaneamente \\
Nel 1900 viene chiamato alla cattedra di filosofia non so dove \\
Nel 900 pubblica "Il riso" \t qua dice che cominco è una categoria umana \\
Nel 1903 pubblica anche "Introduzione alla metafisica" \t che è una sintesi di tutte le sue idee, nel frattempo scrive (1907) "L'evoluzione creatrice" = opera + imp. e sistematica, e opera con + peso teoretico \\
Nel 28 riceve premio nobel per la letteratura \t e nel 32 pubblica "Le due fonti della morale della religione" \\
Bergson è di origine ebraica, e muore nel 41 \t si era avvicinato al cattolicesimo nella fine della sua vita \t pensava che fosse il completamento del giudaeismo \\
Pero non si converte, vuole rimanere ebreo perche clima antisemita \\
Quando i nazisti entrano a parigi, tutti gli ebrei dovevano presentarsi per la schedatura \t ma lui viene esonerato, pero lui va lo stesso \t e muore nel 41 nella parigi occupata dai nazisti

\v

Lo scopo della sua filosofia è difendere la creativita dell uomo, e l'irriducibilita della coscienza a qualsiasi altra cosa \t coscienza o spirito sono a se stanti, non si possono ridurre (contro riduzionismo positivista) \\
Lui non minimizza la dimensione della corporeita pero, e non attaca il mondo materiale \t per capire mondo della coscienza, dobbiamo intedere i risultati della scienza, perche sta nel mondo \\
Non possiamo allontanare la dimensione fisica, perche la coscienza vive nel mondo, e addirittura in un corpo

\ss[Tesi di dottorato]
La sua riflessione parte dalla tesi di dottorato \\
Lui parte leggendo Spencer \t voleva dare un fondamento all'opera di spencer, in particolare a "Primi principi" \\
Indagando Spencer, pero si rende conto che c'è un punto debole \t e ne specula qua \\
Alla meccanica sfugge infatti il tempo dell'esperienza concreta \t nel saggio mette in contrapposizoine il tempo della meccanica, e il tempo della vita \t che chiama durata = tempo dell'esperienza concreta \\
Il tempo della meccanica è un tempo estraneo alla coscienza, perhce è composto da istanti uno esterno all'altro \t sono uno dopo l altro \\
Il tempo della meccanica è anche reversibile \t posso ripetere un esperimento \t inoltre questo tempo è sempre uguale: non c'è un istante che sia superiore e migliore rispetto a un altro \\
Lo definisce anche tempo spazializzato \t perche misurare questo tempo coincide col misurare lo spostamento delle lancette \\
Il tempo della coscienza è pero diverso \t è durata \t è un tempo che vivo con la memoria del passato e l'anticipazione del futuro \\
Quando vivo il presente è il risultato di tutto cio che è stato, e allo stesso tempo posso progettare in avanti, posso anticipare \\
Fuori dalla coscienza il passato non è piu e il futuro non è ancora \t un ente non porta con se il suo passato etc. \\
Inoltre tempo della coscienza non è fatto di istanti uguali, ma diversi e non lineari \t la coscienza è come un gomitolo di lana, in cui ci sono momenti + importanti di altri \\
La percezione che ho del tempo è diversa dal tempo reale, e non è neanche un tempo reversibile

\v

La durata vissuta non è il tempo spazializzato della meccanica \t c'è differenza qualitativa \\
Quando parlo di tempo non posso parlare solo di quello dlela meccanica, ma devo trattare anche il tempo della coscienza \t che è il tempo della vita, ed è sempre nuovo \t non posso prevedere cosa succede \\
Sono due tempi diversi con finalita diverse \t il tempo della meccanica va benissimo per le finalita pratica della scienza, che deve costruire teorie valide \\
Ma questa dimensione è inadeguata se la applichiamo alla coscienza \t non possiamo analizzarla con la scienza \t bisogna tenere conto dei diversi aspetti della realta, che vanno analizzati diversamente \\
Positivismo vuole analizzare tutti i fatti con i mezzi della scienza \t questo ne rappresenta il fallimento, dalla presuzione

\v

L'idea di durata, che considera la caratteristica essenziale della coscienza \t si lega alla liberta e alla difesa della liberta, criticando ogni forma di determinismo \\
Critica il determinismo quando crede di poter indagare la coscienza \\
Sia i sostenitori del determinsmo, ma anche i sostenitori del libero arbitrio possono essere considerati entrambi determinsti \\
L'errore comune sta nel tenativo di applicare alla coscienza categorie che non le appartengono, e in particolare la categoria causa-effetto \\
Questo perche la coscienza è creativa e imprevedebile \t quindi non posso leggerla secondo la causalita, che contraddirebbe la natura della coscienza \\
\bi
\item i deterministi vanno a cercare le cause di un azione, che ricerco in un orizzonte di necessita \t non c
\item sost. libero arbitrio dicono che c'è sempre la stessa causa alla radice dell'azione libera, che è la volonta \t lo schema è uguale: voglio quindi sono libero
\ei
La mia azione è o determinata da una causa necessaria, o dalla volonta che è sempre causa della mia azione \t quindi entrambi leggono la coscienza come atti distinti \\
Non c'è niente di lineare nella coscienza pero \t questo schema fallisce nel tempo, e anche nella liberta \\
Un azione è libera quando dipende da noi \t e dipendere da me vuol dire non dipendere da un atto di volonta, ma quando è totalmente coerente con cio che sono diventato \t quando è coerento con cosa sono in quel memento, con la mia storia \\
Siamo liberi quando la mia azione non nasce in modo lineare da causa A ad effetto B, ma nasce da questo gomitolo da cui esce azione \\
QUesto è anche perche davanti a eventi uguali, uomini scelgono diveramente

\v

Noi pero compiano anche molte azioni che non ci appartengono \t e infatti quelle azioni sono abitudini, che è prevedibile, mecanica \\
L'abitudine è prevedibile \t è un orizzonte lineare \t non sono negative, perche viviamo anche nell abitudine, che pero non è liberta \t non esprime noi stessi, significa stare in superificie \\
Piu vado alla radice meno le mie azioni sono abitudinarie \t piu sono prevedibile \t sono scelte meccaniche, non libere
\end{document}
